\chapter*{Preface}
\addcontentsline{toc}{chapter}{Preface}
\markboth{Preface}{Preface}

\newthought{This manual} is the complete, ground-up reference for \Jarvis---a
\textsc{yaml}-driven framework for orchestrating likelihood-based parameter scans in
high-energy physics. It is a living document, revised in step with each \Jarvis\ release;
this build documents \Jversion. It is written for the physicist who wants to go from an empty
directory to a finished, reproducible scan, and who then wants to understand every knob
along the way.

\section*{Who this is for}

You will get the most from this book if you are comfortable running commands in a Unix
shell and editing \textsc{yaml} files. No prior experience with \Jarvis\ is assumed.
Part~\ref{part:getting-started} takes a newcomer end to end; the later parts are a
reference you can dip into key by key.

\section*{How to read it}

\begin{itemize}
  \item \textbf{New users} should read Part~\ref{part:getting-started}
        (\nameref{part:getting-started}) in order, run the quick-start scan, then skim
        Part~\ref{part:task-card} to understand the task card.
  \item \textbf{Writing a real scan?} Live in Part~\ref{part:task-card}
        (the \textsc{yaml} reference) and Part~\ref{part:backends} (calculators, operas,
        and library dependencies).
  \item \textbf{Choosing or tuning a sampler?} Part~\ref{part:samplers} has one focused
        section per method, plus a decision guide and a master comparison table.
  \item \textbf{Running, monitoring, and post-processing?} Part~\ref{part:running}.
  \item \textbf{Squeezing out performance?} Part~\ref{part:advanced}.
  \item The \textbf{appendices} are lookup tables: every \textsc{cli} flag, every
        \textsc{yaml} key, every distribution, every sampler.
\end{itemize}

\section*{Conventions}

\begin{description}[style=nextline, leftmargin=1.2em]
  \item[\code{monospace}] Code, file names, \textsc{yaml} keys, and shell commands.
  \item[\marker{\&J/}] The path \emph{marker} that Jarvis expands at runtime
        (Chapter~\ref{ch:core-concepts}).
  \item[\token{@SampleID}, \token{@Sdir}, \token{@PackID}] Runtime \emph{tokens} substituted
        per sample.
\end{description}

\noindent Three kinds of call-outs appear throughout:

\begin{jnote}
A clarification or a detail that is easy to miss.
\end{jnote}

\begin{jtip}
A recommended practice that will save you time.
\end{jtip}

\begin{jwarn}
A foot-gun. For example: the calculator key is spelled \yamlkey{make\_paraller}
(legacy spelling, kept on purpose), and a task card uses \emph{either} \yamlkey{Calculators}
\emph{or} \yamlkey{Operas} as its workflow backend---not both.
\end{jwarn}

\section*{What's new in this build}

As of \Jversion, \Jarvis\ ships an optional high-throughput runtime, configured entirely through the
\yamlkey{Runtime} task-card block (Chapter~\ref{ch:runtime-block}). Scans can be proposed and
persisted in batches (\yamlkey{batch\_size}), run across multiple worker processes
(\yamlkey{mode: process}, \yamlkey{Operas}-only for now), and persist through a separate writer
process (\yamlkey{writer: process}); per-sample folders are now created lazily
(\yamlkey{sample\_artifacts: auto}), with failed points still guaranteed a complete log. None of this
changes existing task cards, the \textsc{cli}, the output layout, or checkpoint behaviour---with no
\yamlkey{Runtime} block the engine runs exactly as before. Measure any of it reproducibly with
\cli{-{}-benchmark} (Chapter~\ref{ch:performance}).

\section*{Source of truth}

\Jarvis\ evolves quickly. Where this manual and the running code ever disagree, the code
wins; please report the gap on the issue tracker. Every \textsc{yaml} key printed here was
checked against the runtime configuration loader and the per-sampler \textsc{json} schemas
shipped in \schemaref{}.

\vspace{1\baselineskip}
\noindent If \Jarvis\ saves you a week of bookkeeping, it has done its job.
